% Hypergeometric-distribution.tex
% Topic: Hypergeometric Distribution
% (Advanced topic — outside the HSC Mathematics Extension 2 syllabus)

The \textbf{hypergeometric distribution} models the number of successes in a sample drawn \emph{without replacement} from a finite population.

\textbf{Setup.} A population of $N$ items contains $K$ successes (and $N-K$ failures). A sample of $n$ items is drawn without replacement. Let $X$ be the number of successes in the sample.

\textbf{Probability formula:}
\[
P(X = k) = \frac{\dbinom{K}{k}\dbinom{N-K}{n-k}}{\dbinom{N}{n}}, \qquad \max(0,\,n-(N-K)) \leq k \leq \min(n,K).
\]

\textbf{Mean and variance:}
\[
E(X) = \frac{nK}{N}, \qquad \mathrm{Var}(X) = \frac{nK(N-K)(N-n)}{N^2(N-1)}.
\]

\textbf{Contrast with binomial.} The binomial distribution models sampling \emph{with replacement}; the hypergeometric distribution models sampling \emph{without replacement}. For large $N$, the two distributions are approximately the same.

\begin{remark}
The hypergeometric distribution arises naturally in quality control, card problems, and ecological sampling.
\end{remark}

% Placeholder: examples to be added.
